%% kapa: exam=1 round=28 year=2017 period=1 form=A subject=real_estate subject_ko=부동산학원론 seq=5 number=85
%% source: Q-Net 감정평가사 기출문제 (cst00309) artlSeq=5199603 file=extracted/1008/5199603/2017년 제28회 감정평가사 1차시험 문제지/1교시 A형.pdf
%% conversion: kordoc 4.24.51 -> tools/gen1.py (auto); review=proofread-sonnet; flags=none
\begin{question}{85}
주택금융에 관한 설명으로 옳은 것을 모두 고른 것은?
\begin{framedtext}
ㄱ. 주택금융은 주택수요자에게 자금을 융자해 줌으로써 주택구매력을 높여준다. \par
ㄴ. 주택소비금융은 주택을 구입하려는 사람이 신용을 담보로 제공하고, 자금을 제공받는 형태의 금융을 의미한다. \par
ㄷ. 주택개발금융은 서민에게 주택을 담보로 하고 자금을 융자해주는 실수요자금융이다. \par
ㄹ. 주택자금융자는 주로 장기융자 형태이므로, 대출기관의 유동성 제약이 발생할 우려가 있어 주택저당채권의 유동화 필요성이 있다.
\end{framedtext}
\begin{choices}
  \item ㄱ, ㄴ
  \item ㄱ, ㄷ
  \item ㄱ, ㄹ
  \item ㄴ, ㄹ
  \item ㄷ, ㄹ
\end{choices}
\end{question}
\answer{3}
